% use \chaptermark{}
% to alter or adjust the chapter heading in the running head


%%%%%%%%%%%%%%%%%%%%%%%%%%%%%%%%%%%%%%%%%%%%%%%%%%%%%%%%%%
%%%
%%%  A Semigroup Approach to the Road Coloring Problem
%%%
%%%%%%%%%%%%%%%%%%%%%%%%%%%%%%%%%%%%%%%%%%%%%%%%%%%%%%%%%%
\section{Introduction}  
\label{sec:4.1}

Suppose one is given a strongly connected digraph (that is, there is a directed path between any two vertices).  
One can think of the vertices of the graph as buildings, for example, and the directed arcs as unnamed, one-way roads connecting the buildings. 
Suppose, further, that in each building there is a person. 
Under what conditions can one color the roads, so that the same set of instructions allows each person to get to the same building.

This problem, whose precise formulation is given below, was first posed by
Adler, Goodwyn, and Weiss \cite{AGW} in a paper more than twenty years ago. 
Due to the informal interpretation given above, it has become known as the road coloring problem, and the set of directions which gets everyone to the same building is known as the synchronizing instruction.  

Various partial results exist (as our list of references indicates). 
The techniques employed by these papers vary considerably, but almost all have a heavy combinatorial flavor. 
The main contribution of this paper is to translate the problem into semigroup theory and to use the fundamental theorem of finite semigroups to both generalize the problem and to simplify the proofs of many of the past results. 
This unification of technique permits the ability to see many diverse results from a common perspective. In particular, the structure theory of the minimal ideal (kernel) of a finite semigroup plays an essential role. 
We feel that the kernel will eventually play an important role in the complete solution of the problem. 
All of the semigroup theory used is elementary and can be found, for example, in Clifford and Preston \cite{CP}.

The structure of this paper will be as follows. 
After this introduction, we will provide the necessary semigroup theory to state the problem precisely. 
Next we will include a section that shows how the structure of the kernel allows us to translate certain graph theoretic properties  into the language of semigroups. 
In addition, we will generalize the problem in this semigroup context. 
Finally, we will present new semigroup proofs of several of the most well known results.
Examples to illustrate various key ideas will be used throughout the paper.
